
The methodology follows from the observation that MCP tool-calling traces encode researcher reasoning in natural language. If traces contain both explicit structure (tool names, parameter types, result schemas) and implicit reasoning (assumptions in query text, evidence in result text), then a three-layer analysis framework should detect semantic failures invisible to syntactic validation.

\textbf{Three-Layer Architecture.} Layer 1 (Syntactic Structure Validation) checks that MCP traces contain complete records with all required fields (tool name, parameters, result) and validates JSON schemas using standard validation libraries. This layer answers: ''Are traces structurally complete?'' A threshold of 95\% completeness was tested to establish whether traces provide sufficient data for semantic analysis.

Layer 2 (Semantic Query-Parameter Analysis) extracts assumptions from query text using LLM-based NLP. Given a tool call with parameters containing query text such as ''search for pruning techniques reducing effective rank,'' Layer 2 identifies implicit assumptions. Zero-shot prompting with pre-trained LLMs (Claude Sonnet 4.5) and multi-vote consensus (3 independent extractions, $\geq 2/3$ agreement threshold) mitigates hallucination. A threshold of 90\% natural language presence ($\geq 10$ words) was tested to validate that tool calls contain sufficient textual content for extraction.

Layer 3 (Semantic Result-Content Analysis) extracts evidence from result text using the same LLM-based approach. Cross-phase comparison of assumptions extracted from early-phase queries (Phase 1-3) against claims extracted from later-phase results (Phase 4-6) detects semantic contradictions. Constraint inference was tested using sentence-transformer embeddings (all-MiniLM-L6-v2 model, 384-dimensional) with cosine similarity threshold <0.3 to flag contradictions.

\textbf{Design Decision: Zero-Shot LLM Extraction.} Pre-trained LLMs were used without fine-tuning for three reasons. First, the framework must operate without labeled data to satisfy feasibility requirements. Second, pre-trained models like Claude Sonnet 4.5 are trained on scientific literature, providing strong zero-shot capability on research domain text. Third, multi-vote consensus (running three independent extractions and accepting items appearing in $\geq 2/3$ votes) reduces false positives from single-call errors.

\textbf{Design Decision: Sentence-Transformer Embeddings.} For constraint inference, semantic similarity via sentence-transformers was initially chosen based on the hypothesis that contradictory statements discussing the same entity with opposite conclusions would exhibit low similarity. This approach failed (0\% recall), revealing a fundamental mismatch: sentence embeddings optimize for semantic relatedness (topic similarity, paraphrase detection), not logical contradiction. Contradictory statements like ''X decreases'' and ''X increased by 6.02\%'' have high semantic similarity because they discuss the same concept with shared terminology. The negative result establishes a boundary condition for semantic similarity approaches.

\textbf{Data Collection.} The dataset comprises 20 MCP execution traces from research pipeline runs (10 successful, 10 failed), containing 596 tool calls. Each trace is a JSONL file where each line represents one tool call with fields: tool\_name, parameters (including query text), result (including returned content), timestamp, and phase. The parameters and result fields contain natural language encoding assumptions and evidence respectively.

\textbf{Implementation.} All code is implemented in Python 3.10 with modular architecture matching the three-layer design. Layer 1 uses json and pathlib for parsing, regex for word counting. Layers 2-3 use Anthropic API (Claude Sonnet 4.5) for extraction with temperature=0.0 and n\_votes=3. Layer 3 uses sentence-transformers library for all-MiniLM-L6-v2 embeddings and torch for cosine similarity computation. Human annotation validation (Layer 2) uses Cohen's Kappa from scikit-learn for inter-rater agreement.
